% !TeX root = ../papers/nc-p3.tex
\section{大きい五つの台に共通する四次障害と階数十の正則曲線}
\label{sec:horizontal-rectangle-low-fibers}
\begin{theorem}\label{thm:horizontal-rectangle-low-fibers}
112型の水平代表台143,144,145,146,147で、各始点の長さ4が標準値35を持つための必要十分条件は
$b+c\ne0$ である。そのとき全始点の長さ5も標準値56を持つ。
\end{theorem}
\begin{proof}
$b+c=0$ は命題\ref{prop:parallel-antidiagonal-obstruction}で排除した。
残りでは前節までの $u,v,z,w$ を使う。指定した五代表で共通に可逆なのは
\[
 M=uv(u+v)(1+uv)(zw-1)\prod_{q\in\{z,w\}}(q+v)(q-u)(uq-1)(vq+1)
\]
である。以下の論証ではこれ以外の厚い係数の非零を使わない。
最初に $b\ne1$ とし $t,b,x=a,y$ を使って $c=(t^2-b)/(1-b)$ を消去する。
各位置各色の関係フレームには単元小行列式がある。五次までの経路を作ったときに必要な追加の因子は
\[
\begin{split}
 L={}&b^2t^2-b^2xy+b^2x+b^2y^2-b^2+b t^2xy-b t^2x-2bt^2y\\
 &{}-by^2+by+t^2y
\end{split}
\]
のみで、この開集合上では全始点で求める次元が得られる。

$J$ の四つの冪と、それぞれへの $I$ の合成を使う。必要な係数変更は、一般テンソルで
全係数に成立する対角倍率と置換であり、薄い係数と規格化 $B$ の倍率が可逆な範囲で厚い零成分にも特殊化できる。
この点を具体的に確認するには、動かした表示の $B$ がいずれの薄いゲージ規格化の後でも
$x_1,x_4,x_{13},x_{15}$ のいずれかの可逆な基準係数またはその対角倍率倍になることを使う。
実際、各変換の逆も同じ表示式で定義され、窓の行列式に対応する $(zw-1)$ は可逆である。
変換先の $b'\ne1$ は $u'^2\ne1$ と同値である。
$u^2,v^2\ne1$ の開集合では八個の $L$ の像を取り、$M(u^2-1)(v^2-1)$ で飽和すると単位イデアルである。
$u=1,-1$ では $v^2\ne1$ の範囲で $b'\ne1$ の四表示だけを使い、両方で再び単位イデアルを得る。
$v=\pm1,u^2\ne1$ は回転でここに移る。残る $u^2=v^2=1$ では $u=v$ であり、
特別成分で $a,y\ne0,1$ の既知の正則点である。
\path{computations/verify_second_rectangle_low_fibers.py} が追加の一因子と三つの飽和を再現する。
\end{proof}

\begin{proposition}
76型の水平台85で $u=1,z=-1$ の全点は標準Koszul左右AS正則である。
その回転、コサイクル捻りにより $(v,z)=(1,1)$, $(u,z)=(-1,1)$, $(v,z)=(-1,-1)$ の各曲線も全点で正則である。
\end{proposition}
\begin{proof}
最初の曲線のパラメータは $v\ne0,\pm1$ である。
各位置を $(0,2,2,0)$ だけ平行移動した座標で、文字0の四次中心核の自由列の最初と最後、
文字1の両自由列を使う。これら四式を割った経路の長さ13までの次元は一般に
$(1+T+T^2+T^3)^4$ である。各始点で掃き出しに要る既約因子は順に $196,183,185,175$ 個である。
$v\mapsto1/v$ で移した二位置先のリストとの共通因子を取ると、始点0では14個、次数和31、
始点1では19個、次数和38となる。和集合は26因子、最大次数4である。
各因子の数体について、定義域を保つ一次剰余根で経路と四次中心の次元、四中心式の商の
長さ13までの全四始点の上界を検算する。中心の持上げと一般点のHilbert下界で標数零の全共役点に結論が従う。
有限個の始点に共通する一般の四式を取ることで自由性と左右正則性が得られる。
\path{computations/verify_first_rectangle_rank_ten_curve.py} は元の一変数体上の掃き出しと全26数体を検算する。
\end{proof}
